\subsection{\problemblock{*Reset} Data Block}\label{sec:block-reset}

The \problemblock{*reset} data block is very similar to the
\problemblock{*increment} data block described in \cref{sec:block-increment}.
However, instead of incrementing time or state variables, the
\problemblock{*reset} block sets them equal to input values. All of the input
restrictions discussed for \problemblock{*increment} apply to
\problemblock{*reset}. The variables that can be reset are listed in
\cref{tab:increment-variables}.

A typical use of the block is to reset \statevar{tmark}:

\begin{lstlisting}[style=taosmanual]
*reset tmark=0
\end{lstlisting}

This marks an event during the trajectory. Rocket-motor thrust in propulsion tables
is often made a function of \statevar{tmark} rather than absolute
\statevar{time} or segment time \statevar{tseg}. This allows the motor burn to
extend across several segments and permits motor ignition at any point in the
trajectory. The value of \statevar{tmark} need only be reset to zero at motor
ignition. Other similar variables are \statevar{grmark}, \statevar{plmark}, and
\statevar{fuel}.

Another common use is to set the weight to a known value after missile staging. For
example,

\begin{lstlisting}[style=taosmanual]
*reset wt=1100
\end{lstlisting}

sets the weight to 1100~lb regardless of what has happened previously in the
trajectory.
